%% threat_model.tex

\section{Threat Model and Problem Formalization}
\label{sec:threat-model}

\subsection{System Model}

We consider a multi-agent LLM operating system (exemplified by \system{}) with the following architecture:
\begin{itemize}[leftmargin=*]
    \item \textbf{Agents} are scheduled as processes with per-agent \emph{roles} (e.g., \code{Code\_Reviewer}, \code{Python\_Coder}) drawn from a whitelisted blueprint registry.
    \item \textbf{Tools} (file I/O, shell, web fetch) are syscalls gated by a \code{PermissionChecker} that enforces role-based access control and tenant ownership on the forward execution path.
    \item \textbf{Recovery orchestrator} classifies failures and dispatches a layered chain of recovery strategies organized into 11 conceptual tiers (from lightweight in-place retry up to circuit-breaker human escalation), realized by 14 concrete strategy classes---including reflection injection, \emph{topology mutation}, model fallback, and human escalation. The 11/14 distinction is intentional: the tiers describe \emph{what kind} of failure is being handled (and at what escalation level), while the 14 classes are the concrete implementations, with some tiers containing multiple strategy variants.
    \item \textbf{Semantic core dump} captures an agent's last input, context, and error trace upon crash, for post-mortem diagnosis by an LLM diagnostic module.
\end{itemize}

\subsection{The Privilege Containment Gap at the Recovery Sink}

The system maintains two privilege boundaries (Figure~\ref{fig:trust-boundaries}):
\begin{enumerate}[leftmargin=*]
    \item \textbf{Forward-path boundary}: User input and external tool outputs are treated as untrusted. The \code{PermissionChecker} enforces role-based access control before any tool executes. A \code{Code\_Reviewer} agent, for example, cannot invoke \code{bash} unless its role profile permits it.
    \item \textbf{Recovery-path boundary}: When an agent fails, the recovery orchestrator may dispatch a \code{TopologyMutationStrategy} that reads the semantic core dump, feeds it to a diagnostic LLM, extracts a \code{suggested\_role} from the LLM's JSON response, and passes this role to the node-replacement primitive (\code{resumeNode}). \emph{This boundary is the focus of this paper.} In the baseline configuration, the role-replacement sink performs no authorization check on the value it consumes: the LLM-diagnosed role reaches \code{resumeNode} directly, bypassing the \code{PermissionChecker} that gates the forward path.
\end{enumerate}

\begin{figure}[htbp]
\centering
\begin{tikzpicture}[
    node distance=0.6cm and 0.7cm,
    scale=0.82, transform shape,
    box/.style={draw, rounded corners=2pt, minimum height=0.7cm, minimum width=2.0cm, align=center, font=\small},
    benign/.style={box, fill=green!15, draw=green!60!black},
    attack/.style={box, fill=red!15, draw=red!70!black},
    critical/.style={box, fill=red!70, draw=red!70!black, text=white, font=\small\bfseries},
    arr/.style={-Stealth, thick},
    barr/.style={-Stealth, thick, green!60!black},
    marr/.style={-Stealth, thick, red!70!black, dashed},
    region/.style={draw, dashed, rounded corners=4pt, inner sep=0.35cm}
]
% Benign flow
\node[benign] (user) {Legitimate User};
\node[benign, right=of user] (web1) {Normal Webpage};
\node[benign, right=of web1] (agent1) {Agent Executes};
\node[benign, right=of agent1] (done) {Task Done};
\draw[barr] (user) -- (web1);
\draw[barr] (web1) -- (agent1);
\draw[barr] (agent1) -- (done);

% Malicious flow
\node[attack, below=1.8cm of user] (atk) {Attacker\\(External Provider)};
\node[attack, right=of atk] (mal) {Malicious Webpage\\/ Poisoned API};
\node[attack, right=of mal] (fail) {Agent Crash\\(induced)};
\node[critical, right=of fail] (dump) {Poisoned Core Dump\\(embeds role hint)};
\draw[marr] (atk) -- (mal);
\draw[marr] (mal) -- (fail);
\draw[marr] (fail) -- (dump);

% Recovery path
\node[critical, right=of dump] (orch) {Diagnostic LLM\\$\rightarrow$ suggested\_role};
\node[critical, right=of orch] (esc) {\code{resumeNode}\\(no authz check)};
\draw[marr] (dump) -- (orch);
\draw[marr] (orch) -- (esc);

% Region backgrounds
\begin{scope}[on background layer]
\node[region, fill=green!5, fit=(user)(web1)(agent1)(atk)(mal)(fail), label={[font=\footnotesize\itshape, green!60!black]above:Forward Path (permission-gated)}] {};
\node[region, fill=red!5, fit=(dump)(orch)(esc), label={[font=\footnotesize\itshape, red!70!black]above:Recovery Path (no authorization)}] {};
\end{scope}
\end{tikzpicture}
\caption{Privilege boundaries and the topology-mutation attack chain. The forward path (top, green) enforces role-based access control via the \code{PermissionChecker}; the recovery path (bottom, red) consumes an LLM-diagnosed role at the replacement sink without re-checking authorization, enabling privilege escalation.}
\label{fig:trust-boundaries}
\end{figure}

\subsection{Attacker Model}
\label{sec:attacker-model}

We assume the attacker is an \textbf{external content provider}---e.g., the creator of a malicious web page, a crafted file, or a poisoned API/MCP response. Concretely, the attacker controls \emph{external content} that the agent processes during normal operation:
\begin{itemize}[leftmargin=*]
    \item Web pages fetched via \code{web\_fetch} or \code{web\_scrape}
    \item Files read from the filesystem (if accessible)
    \item Tool outputs from external integrations (MCP servers, third-party APIs)
\end{itemize}

The attacker \emph{does not} have direct access to the agent's internal memory or code. In particular, the attacker cannot:
\begin{itemize}[leftmargin=*]
    \item Modify the agent's role, permission profile, or whitelisted blueprint
    \item Alter the LLM model, system prompt, or model provider
    \item Tamper with the recovery orchestrator's strategy dispatch logic or the \code{PermissionChecker}
    \item Inject code into the trusted computing base (kernel, runtime, or recovery machinery)
\end{itemize}

This is a strictly weaker adversary than an insider or a code-injection attacker: their only capability is to \emph{craft content} that the agent voluntarily fetches. The threat, therefore, must exploit a path by which such externally-fetched content \emph{acquires} the authority to trigger a privileged role replacement without ever touching the trusted computing base directly. The recovery path's topology-mutation strategy provides exactly such a path.

\subsection{Attack Chain: Topology-Mutation Privilege Escalation}
\label{sec:attack-chain}

The attack follows a five-stage chain (Figure~\ref{fig:attack-chain}). The first three stages occur on the \emph{forward} path and are indistinguishable from benign operation; the last two occur on the \emph{recovery} path and constitute the privilege escalation.

\begin{enumerate}[label=(\roman*),leftmargin=*]
    \item \textbf{Content Delivery.} Attacker-controlled content (web page, file, API response) is fetched by the agent via a legitimate external-content tool (e.g., \code{web\_fetch}). At this stage the content is correctly treated as untrusted by the forward-path boundary.
    \item \textbf{Failure Induction.} The content is crafted to trigger an agent crash---a capability mismatch, a verification failure, or a semantic crash that the recovery orchestrator routinely handles. No exotic bug is required; the orchestrator is designed to recover from precisely these failures via topology mutation.
    \item \textbf{Payload Embedding.} The adversarial payload---a role-escalation hint such as \code{"suggested\_role": "System\_Admin"}---is embedded in the semantic core dump that captures the agent's last input and context on crash. Because the capture site records the dump as an opaque string, the payload is preserved verbatim.
    \item \textbf{Diagnostic LLM Consultation.} The \code{TopologyMutationStrategy} feeds the core dump to a diagnostic LLM and asks it to recommend a replacement role. The LLM, reading the embedded hint, emits the attacker-suggested high-privilege role in its JSON response.
    \item \textbf{Unauthorized Role Replacement.} The extracted \code{suggested\_role} reaches \code{resumeNode}---the node-replacement primitive---without an authorization check. The agent is replaced with the attacker-suggested role, escalating privileges.
\end{enumerate}

\begin{figure}[htbp]
\centering
\begin{tikzpicture}[
    node distance=0.6cm and 0.5cm,
    stage/.style={draw, rounded corners=2pt, minimum height=0.8cm, minimum width=2.2cm, align=center, font=\footnotesize},
    fwd/.style={stage, fill=green!12, draw=green!60!black},
    rev/.style={stage, fill=red!15, draw=red!70!black},
    final/.style={stage, fill=red!70, draw=red!70!black, text=white, font=\footnotesize\bfseries},
    arr/.style={-Stealth, thick, gray!70},
    lbl/.style={font=\tiny\itshape, gray!60!black, align=center}
]
% Stage chain
\node[fwd] (s1) {(i) Content\\Delivery};
\node[fwd, right=of s1] (s2) {(ii) Failure\\Induction};
\node[fwd, right=of s2] (s3) {(iii) Payload\\Embedding};
\node[rev, right=of s3] (s4) {(iv) Diagnostic\\LLM Consult};
\node[final, right=of s4] (s5) {(v) Unauthorized\\Role Replace};

\draw[arr] (s1) -- (s2);
\draw[arr] (s2) -- (s3);
\draw[arr] (s3) -- (s4);
\draw[arr] (s4) -- (s5);

% Stage descriptions
\node[lbl, below=0.15cm of s1] {web\_fetch / file\_read};
\node[lbl, below=0.15cm of s2] {capability mismatch};
\node[lbl, below=0.15cm of s3] {core dump + role hint};
\node[lbl, below=0.15cm of s4] {LLM emits suggested\_role};
\node[lbl, below=0.15cm of s5] {\code{resumeNode} (no authz)};

% Boundary marker
\draw[dashed, red!50, thick] ($(s3.east)+(0.25,-0.8)$) -- ($(s3.east)+(0.25,0.8)$);
\node[font=\tiny\itshape, red!50, right=0.05cm of s3.east, yshift=0.9cm] {privilege boundary crossed};

% Region labels
\node[font=\footnotesize\itshape, green!60!black] at ($(s2.north)+(0,0.5)$) {Forward Path (benign-looking)};
\node[font=\footnotesize\itshape, red!70!black] at ($(s4.north)+(0.5,0.5)$) {Recovery Path (escalation)};
\end{tikzpicture}
\caption{The five-stage topology-mutation privilege-escalation chain. Stages (i)--(iii) occur on the forward path and are indistinguishable from benign operation; stages (iv)--(v) occur on the recovery path and constitute the escalation. The privilege boundary is crossed between stages (iii) and (iv), where the recovery pipeline consumes an LLM-diagnosed role without re-checking authorization.}
\label{fig:attack-chain}
\end{figure}

The critical property of this chain is that \emph{no single stage is anomalous}: each stage is a feature of the self-healing design, and the attacker merely composes them. The topology-mutation strategy is \emph{designed} to consult a diagnostic LLM and act on its recommendation; the bug is that ``acting on'' the recommendation bypasses the authorization check that the forward path enforces.

\subsection{Formalization: Two Complementary Framings}

We formalize the privilege containment gap along two complementary dimensions, both of which motivate the same structural fix.

\subsubsection{CWE-862 (Missing Authorization) at the Recovery Sink}
The role-replacement primitive \code{resumeNode(nodeId, report, workflowId)} consumes a \code{suggested\_role} field extracted from the diagnostic LLM's JSON response. The role value is \emph{influenced by content that originated outside the trusted computing base} (the core dump embeds external content the agent was processing when it crashed), yet the sink performs no authorization check on this value~\cite{cwe862}. This is a textbook instance of CWE-862: a program performs a privileged operation (role replacement) without verifying that the caller is authorized to request that operation.

Formally, let $R_{\text{old}}$ be the failing agent's role and $R_{\text{new}}$ be the LLM-diagnosed replacement role. The baseline sink executes the replacement whenever $R_{\text{new}} \neq \bot$, with no check that $R_{\text{new}}$ is in the legitimate downgrade set of $R_{\text{old}}$. The attacker's goal is to induce $R_{\text{new}}$ such that $P(R_{\text{new}}) > P(R_{\text{old}})$, where $P(\cdot)$ denotes the privilege level.

\subsubsection{Confused Deputy at the Orchestrator}
The recovery orchestrator is a \emph{privileged deputy}: it holds the capability to swap agent roles---a capability not granted to ordinary tool calls on the forward path. In the baseline configuration, the orchestrator exercises this capability on the basis of a diagnostic LLM's recommendation, whose input (the core dump) was supplied by attacker-influenced content. This is the Confused Deputy problem~\cite{hardy1988confused}: a privileged program is tricked by a less-privileged caller into misusing its authority on the caller's behalf.

The classic Confused Deputy fix is \emph{capability-based authority}: the deputy must present a capability that justifies the action, not merely act on behalf of a request. In our setting, this means the orchestrator must verify that the proposed role replacement is \emph{within the failing agent's authority to request}---i.e., that the new role does not exceed the old role's privileges---regardless of what the diagnostic LLM suggested.

\subsection{Why Per-Strategy Checks Are Insufficient}
\label{sec:why-centralize}

A natural response to CWE-862 at the recovery sink is to add an authorization check \emph{inside} each side-effecting recovery strategy: \code{TopologyMutationStrategy} checks before replacing a role, \code{ModelFallbackStrategy} checks before switching models, and so on. This decentralized approach suffers from two failure modes that are well-known in the cross-cutting concern literature:

\begin{enumerate}[leftmargin=*]
    \item \textbf{Gaps.} A strategy author who forgets or mis-scopes a check (as in the baseline \code{parseAndValidate(validate=false)} call) leaves a silent hole, because no other component re-checks the action. The baseline vulnerability in \system{} is exactly such a gap: the \code{validate=false} flag was a deliberate experimental switch, but in a production codebase, an equivalent gap could arise from a forgotten check in a newly added strategy.
    \item \textbf{Inconsistency.} Different strategies implement slightly different policies, so the effective privilege boundary depends on which strategy the orchestrator happens to dispatch. A role replacement via \code{TopologyMutationStrategy} might be checked, while a role replacement via a future \code{RolePromotionStrategy} might not be.
\end{enumerate}

The principled fix is to centralize the concern at a single choke point---the \reauthgate{} (\S\ref{sec:defense}). This localization is itself an OS principle: classical kernels enforce privilege invariants at well-defined boundaries (syscall entry, exception dispatch), not inside each driver or subsystem.
